\documentclass[11pt]{article}

\usepackage[hmargin=0.78in,vmargin=0.68in]{geometry}
\usepackage[T1]{fontenc}
\usepackage{lmodern}
\usepackage{microtype}
\usepackage{amsmath,amssymb,graphicx}
\usepackage[dvipsnames]{xcolor}
\usepackage{booktabs,tabularx,array}
\usepackage{enumitem}
\usepackage{titlesec}
\usepackage{caption}
\usepackage[hidelinks]{hyperref}

\setlength{\parindent}{0pt}
\setlength{\parskip}{4pt}
\setlist[enumerate]{leftmargin=1.45em,itemsep=2pt,topsep=3pt}
\setlist[itemize]{leftmargin=1.45em,itemsep=2pt,topsep=3pt}
\titlespacing*{\section}{0pt}{9pt}{3pt}
\titlespacing*{\subsection}{0pt}{7pt}{2pt}
\titleformat{\section}{\large\bfseries}{}{0pt}{}
\titleformat{\subsection}{\normalsize\bfseries}{}{0pt}{}
\captionsetup{font=small,labelfont=bf,skip=4pt}
\newcolumntype{Y}{>{\raggedright\arraybackslash}X}
\newcommand{\yes}{\textcolor{ForestGreen}{yes}}
\newcommand{\no}{\textcolor{BrickRed}{no}}

\begin{document}

\begin{center}
{\LARGE\bfseries Reproductive Closure in the Full Lifecycle Model}\par
\vspace{3pt}
{\large Numerical audit and transition requirements}\par
\vspace{3pt}
{\small August 11, 2026}
\end{center}

\vspace{-2pt}
This memo carries out the three checks listed at the end of the companion note
\emph{Fertility, Population, and House Prices}: baseline reproduction, the
shape of the reproductive and housing-scale mappings, and the requirements for
a calendar-time transition.  The conclusion is not that the alternative
closure is conceptually wrong.  The conclusion is that it is \emph{not yet a
closure of the maintained quantitative model}.

\section*{1. Bottom line}

The proposed equilibrium uses two conditions.  Effective reproduction selects
the housing cost, and housing clearing then selects population scale:
\begin{equation}
 B(r;\Theta)=E(r;\Theta),
 \qquad
 S=\frac{H^S(r;\Theta)}{D(r;\Theta)}.
 \label{eq:closure}
\end{equation}
Here $B$ is the flow of children who mature into new adult households, $E$ is
the entrant-household flow needed to reproduce the adult distribution, $D$ is
housing demand per adult household, and $H^S$ is aggregate housing supply.  The
first condition is the missing demographic closure; the second is the existing
housing-market scale condition.

The full-model audit gives four results.
\begin{enumerate}
 \item \textbf{The maintained point is not reproductive.}  At the retained E5b
 parameters, $B/E=0.82784$.  The model produces only 82.8 percent of the entrant
 households needed to reproduce its normalized adult population.
 \item \textbf{Cheaper housing moves fertility in the right direction but is
 quantitatively much too weak.}  Lowering the asset price to one percent of its
 maintained value raises $B/E$ only to $0.83600$.  The residual is monotone over
 the tested grid but never crosses zero.
 \item \textbf{The housing block has the opposite quantitative property.}
 The scale map $H^S/D$ is steeply increasing in the housing cost.  It falls to
 $0.000197$ at one percent of the maintained price.  Thus population scale
 disappears much faster than fertility approaches replacement.
 \item \textbf{A full-model transition has not been solved.}  The stationary
 calculations establish endpoint conditions and reject existence at the
 maintained parameter vector.  They do not establish an impact response,
 monotone adjustment, transition speed, or the claim that the current economy
 is already on such a path.
\end{enumerate}

The practical recommendation is therefore to retain the simple model as a
clean analytical illustration, but not yet to replace the paper's population
closure with equation~\eqref{eq:closure}.  Doing that requires an explicit
demographic re-estimation and, if the paper wants transition language, a new
calendar-time solution method.

\newpage
\section*{2. Source candidate and status check}

The audit uses the retained July 25 E5b parameter vector evaluated under the
corrected August 9 target contract.  Its exact tight preflight has loss
$385.87549$, price $P_0=0.815625$, and housing-market residual
$1.11\times10^{-6}$.  Two exact repeats agree in every target moment.

There is a project-status discrepancy worth recording.  The August 9 production
and collector jobs did finish, although the live status note still described
them as running.  The collector's nominal new winner has loss $386.68899$, which
is worse than the retained incumbent.  Because that collector searched only the
new production chains, it was not incumbent-safe.  The retained E5b point
remains the best known point under the corrected contract.  Neither point is a
promoted empirical calibration while the first-birth rooms-code hold remains
open.  Appendix A reports every target row and every free parameter.

\section*{3. Baseline reproduction accounting}

All flows below are per normalized adult household and per four-year model
period.  The exact accounting is:
\begin{align}
 E_0&=0.06173346,
 &\text{required entrant households},\notag\\
 \text{births}_0&=0.10754231,
 &\text{children born},\notag\\
 B_0&=0.05110564,
 &\text{mature entrant households},\notag\\[-2pt]
 \frac{B_0}{E_0}
 &=
 \underbrace{\frac{\text{births}_0}{E_0}}_{1.74204}
 \times
 \underbrace{0.5}_{\substack{\text{children per}\\\text{new household}}}
 \times
 \underbrace{0.95043}_{\substack{\text{maturation and}\\\text{exit yield}}}
 =0.82784.
 \label{eq:accounting}
\end{align}
The 17.2 percent reproduction shortfall therefore combines completed family
size, the external convention that two mature children form one new household,
and the fact that not every child becomes an entrant before parental exit or
terminal lifecycle truncation.

The reported TFR is not the same object.  E5b reports $1.90357$ because the
top-coded $3+$ parity state receives a weight of $3.60236$.  The actual KFE
transition counts that state as three children.  Reweighting the mature flow by
the reported top-code convention raises the diagnostic ratio to $0.90319$, but
still does not produce replacement.  Even at one percent of the maintained
price, the favorable top-code diagnostic reaches only $0.91326$.

This distinction is substantive rather than cosmetic.  $B/E$ depends on the
child-to-household conversion, child maturation, parental survival, and the
terminal lifecycle convention.  It is internally correct model accounting, but
it is not yet an empirically validated net-reproduction rate.

\section*{4. The two equilibrium mappings}

The first panel of Figure~\ref{fig:maintained} recomputes the normalized
stationary distribution at fifteen fixed prices between $0.01P_0$ and $2P_0$.
The second panel evaluates the population-scale map at the same points.

\begin{figure}[ht!]
\centering
\includegraphics[width=0.98\textwidth]{../../../output/model/closed_reproductive_closure_audit_20260811/maintained_closure_grid.png}
\caption{\textbf{Maintained E5b preferences.}  Effective reproduction falls
with the housing cost, as the proposed mechanism requires, but it remains well
below one throughout the grid.  The top-code reweighting is a convention
sensitivity, not the model's KFE flow.  The housing scale map is increasing and
passes through one at the maintained market-clearing price.}
\label{fig:maintained}
\end{figure}

The local elasticities around the maintained point summarize the mismatch:
\begin{equation}
 \frac{d\log(B/E)}{d\log r}\simeq-0.0107,
 \qquad
 \frac{d\log(H^S/D)}{d\log r}\simeq2.637.
 \label{eq:elasticities}
\end{equation}
The fertility feedback has the correct sign, but the housing-scale response is
roughly two orders of magnitude larger.  Most important, existence fails:
$\mathcal R(r)=B(r)-E(r)$ is negative at every tested price.  Therefore the
maintained parameter vector has no positive closed reproductive steady state
on this grid, and there are no old and new full-model endpoints across which to
decompose population.

\section*{5. What it takes to manufacture two endpoints}

The following exercise is deliberately diagnostic.  First, hold the maintained
price fixed and raise the child-utility parameter until $B/E=1$.  This requires
\begin{equation}
 \psi_{\mathrm{child}}: 0.58502\longrightarrow1.38574.
 \label{eq:psi_normalization}
\end{equation}
The normalization pushes top-coded TFR to $2.35230$ and childlessness to
$0.02379$, compared with targets $1.918$ and $0.188$.  It is therefore not an
innocuous scale normalization; it changes the calibration sharply.

Second, measure the greatest fertility recovery available between the
maintained price and one percent of that price.  Even after imposing
replacement at $P_0$, this extreme price decline raises $B/E$ by only
$0.00373$, or 0.37 percentage points.  The constructed preference shock in
Figure~\ref{fig:constructed} uses one half of that limited capacity: it lowers
$B/E$ at the old price by only $0.00187$.  The new reproductive root nevertheless
requires the price to fall to $0.536P_0$.

\begin{figure}[ht!]
\centering
\includegraphics[width=0.98\textwidth]{../../../output/model/closed_reproductive_closure_audit_20260811/constructed_normalized_closure.png}
\caption{\textbf{Constructed replacement normalization.}  The orange schedule
is shifted down by half of the maximum feedback available over the tested price
range.  Restoring replacement requires a 46.4 percent price decline.  The
housing market then supports only 22.8 percent of the old population scale;
the direct preference-and-composition term is essentially zero in this
particular diagnostic.}
\label{fig:constructed}
\end{figure}

The population-scale decomposition is:
\begin{align}
 \Delta\log S
 &=\underbrace{\log\frac{F(P_1,\Theta_0)}{F(P_0,\Theta_0)}}_{-1.47980}
 +\underbrace{\log\frac{F(P_1,\Theta_1)}{F(P_1,\Theta_0)}}_{0.00000}
 =-1.47980,\notag\\
 \frac{S_1}{S_0}&=0.22768.
 \label{eq:scale_decomposition}
\end{align}
This example verifies the comparative-static sign but exposes the quantitative
fragility: reversing a 0.187 percentage-point reproduction shortfall requires
a 77.2 percent population contraction.  A shock near the maintained 17.2
percent shortfall has no root.

\section*{6. Demographic conventions are part of the model}

Figure~\ref{fig:architecture} evaluates the child-maturation issue in two ways.
Simply switching the retained E5b parameters from the historical shared child
clock to independent child maturation raises $B/E$ to $1.11788$.  But that
same-theta model is no longer calibrated: its TFR jumps to $2.69920$.  The
August 5 independently matured model, after re-estimation, instead has
$B/E=0.78517$.

\begin{figure}[ht!]
\centering
\includegraphics[width=0.78\textwidth]{../../../output/model/closed_reproductive_closure_audit_20260811/architecture_sensitivity.png}
\caption{\textbf{Architecture sensitivity.}  Changing maturation at fixed
parameters alters both fertility policies and mature-child flows.  Re-estimating
the alternative architecture reverses the apparent replacement result.  The
closure cannot be repaired by relabeling one accounting field.}
\label{fig:architecture}
\end{figure}

Three choices must be fixed before a reproductive equilibrium can be estimated:
\begin{enumerate}
 \item whether the $3+$ state produces three children or its empirical
 top-coded mean in the mature-child flow;
 \item whether children share the most recent birth's clock or mature
 independently; and
 \item whether children can mature after the parental household dies.
\end{enumerate}
The maintained code currently uses three children in the $3+$ transition, the
shared clock, and parental survival before child maturation.  Those conventions
should either remain the stated demographic model or be revised and then
re-estimated.  They should not be changed solely to force $B/E=1$.

\clearpage
\section*{7. What a transition interpretation requires}

The stationary closure can compare endpoints only after two reproductive roots
exist.  Interpreting today's low fertility as a transition is a different
exercise.  Let $\mu_t(x)$ be the unnormalized distribution of adult households
over the full beginning-of-period state $x$, and let $Q_t$ carry surviving
adults through saving, tenure, fertility, income, adult aging, and child aging.
If $m_t(x)$ is the flow of children maturing into entrant-household units and
$\nu_E$ is their initial wealth-income distribution, the demographic law is:
\begin{equation}
 B_t=m_t'\mu_t,
 \qquad
 \mu_{t+1}=Q_t'\mu_t+B_t\nu_E.
 \label{eq:calendar_demography}
\end{equation}
Equation~\eqref{eq:calendar_demography} makes the cohort lag explicit.  One
must not impose $B_t=E_t$ along the path: equality holds only at a reproductive
steady state.  Children born today mature many years later, while today's
mature entrants reflect earlier fertility.

Household choices also require time-indexed value functions because saving,
ownership, and fertility depend on anticipated future prices:
\begin{equation}
 V_t(x)=\max_{d_t}\left\{u_t(x,d_t)
 +\beta\,\mathbb E_t[V_{t+1}(x')]\right\}.
 \label{eq:calendar_bellman}
\end{equation}
A sequence of isolated stationary solutions is not equivalent to this problem.

Finally, owned housing is durable.  A literal transition needs a predetermined
stock and a construction, demolition, or vacancy margin:
\begin{equation}
 K_{t+1}=(1-\delta_H)K_t+X_t-Z_t,
 \qquad
 H_t^D+V_t^{\mathrm{vac}}=K_t.
 \label{eq:housing_stock}
\end{equation}
With a gross safe return $R$, the deterministic user cost must include the
next-period asset price:
\begin{equation}
 r_t=(R+\delta_H+\tau_t^p)P_t-P_{t+1}.
 \label{eq:dynamic_user_cost}
\end{equation}
At a constant asset price this reduces to the paper's stationary user-cost
relation.  Along a transition, capital gains or losses matter.

The initial adult distribution and housing stock are predetermined.  The
asset price can jump at the preference change, but population, age composition,
the stock of houses, liquid positions, and owned-house positions cannot.  A
terminal condition must converge to a new reproductive steady state.  A full
implementation also needs date-by-date pension and property-tax accounting,
and it must decide how the model's exogenous entrant-wealth distribution is
reconciled with parental bequests.

The current static supply curve can therefore support stationary endpoint
comparisons, or an explicitly labelled quasi-static approximation in which the
housing stock adjusts instantly.  It cannot support a literal durable-housing
transition.  The statement that durability changes only the path is itself
conditional on choosing an investment technology whose stationary locus
reproduces the existing supply curve.

\section*{8. Recommended use in the paper}

There are two coherent routes.

\textbf{Stationary extension.}  Keep the representative-family model as the
analytical result.  Define $B=E$ as a candidate full-model equilibrium concept,
but state that the current E5b estimate fails the existence check.  Re-estimate
the quantitative model only after fixing the demographic conventions and adding
an empirical net-reproduction object.  Report $B/E$, conventional TFR, completed
family-size composition, and population separately.

\textbf{Transition model.}  Calibrate an old reproductive steady state, treat
the preference change as a permanent shock, and solve equations
\eqref{eq:calendar_demography}--\eqref{eq:dynamic_user_cost}.  The present E5b
stationary distribution cannot be relabelled as a closed-economy transition:
its age distribution is held stationary by the maintained entrant injection
even though its own children replace only 82.8 percent of those entrants.

For the current draft, the first route is safer and already useful.  It lets the
paper explain the feedback clearly while being explicit that the full model has
not yet delivered the two reproductive endpoints or their transition.

\clearpage
\section*{Appendix A. Complete source calibration readout}

The tables below report the exact current-contract preflight at the retained
E5b parameter vector.  The target system is overidentified: twelve target rows
discipline ten estimated parameters, with three additional externally fixed
parameters.  The scalar loss is $385.87549$.  The first-birth rooms row remains
under the active rooms-code hold.

\begin{table}[ht!]
\centering
\caption{Complete target fit}
\scriptsize
\setlength{\tabcolsep}{3.2pt}
\begin{tabularx}{\textwidth}{Yrrrrr}
\toprule
Target moment & Target & Model & Gap & Weight & Loss contribution\\
\midrule
Completed fertility (reported as TFR) & 1.9180 & 1.9036 & $-0.0144$ & 1,425.739 & 0.297\\
Childlessness & 0.1880 & 0.0802 & $-0.1078$ & 17,180.744 & 199.700\\
Mean age at first birth & 25.3106 & 25.6374 & 0.3269 & 16.000 & 1.710\\
Share of first births at age 30+ & 0.2701 & 0.3323 & 0.0623 & 15,625.000 & 60.590\\
Rooms after first birth, event time $+3$ & 0.8049 & 0.6572 & $-0.1477$ & 35.735 & 0.780\\
Prime-age rooms: 3+ minus 1--2 children & 0.3677 & 0.3722 & 0.0045 & 2,958.515 & 0.061\\
Parent--childless ownership gap & 0.1677 & 0.1659 & $-0.0018$ & 14,229.591 & 0.045\\
Aggregate ownership rate & 0.5755 & 0.5638 & $-0.0117$ & 1,207.846 & 0.166\\
Mean occupied rooms, ages 18--85 & 5.7800 & 6.0755 & 0.2955 & 11.973 & 1.046\\
Wealth / annual gross labor earnings & 6.8731 & 5.4548 & $-1.4183$ & 6.288 & 12.649\\
Annual bequest flow / aggregate wealth & 0.0088 & 0.0091 & 0.0003 & 5,165,289.256 & 0.399\\
Old-age estate-wealth p90/p50 & 3.4481 & 2.0684 & $-1.3797$ & 56.960 & 108.433\\
\midrule
\multicolumn{5}{r}{Total loss} & 385.875\\
\bottomrule
\end{tabularx}
\end{table}

\begin{table}[ht!]
\centering
\caption{Complete parameter readout}
\scriptsize
\setlength{\tabcolsep}{4pt}
\begin{tabularx}{\textwidth}{Yrrrrl}
\toprule
Parameter & Estimate & Lower & Upper & Near bound? & Status\\
\midrule
Annual discount factor & 0.991097 & 0.940 & 0.9995 & no & estimated\\
$\Delta\alpha$ & 0.012163 & 0.000 & 0.250 & no & estimated\\
$\Delta\alpha$ at first child & 0.023896 & 0.000 & 0.250 & no & estimated\\
Child utility, $\psi_{\mathrm{child}}$ & 0.585018 & $-3.000$ & 3.000 & no & estimated\\
First-birth logit scale & 3.552006 & 0.020 & 50.000 & no & estimated\\
Continuation-birth logit scale & 1.254006 & 0.020 & 50.000 & no & estimated\\
Bequest curvature, $\chi$ & 1.038033 & 0.100 & 5.000 & no & estimated\\
Housing-supply level, $H_0$ & 8.162391 & 0.200 & 80.000 & no & estimated\\
Bequest shifter, $\theta_0$ & 0.168435 & 0.000 & 8.000 & no & estimated\\
Bequest slope, $\theta_1$ & 0.050046 & 0.020 & 16.000 & \textbf{yes, lower} & estimated\\
Consumption share, $\alpha$ & 0.733 & -- & -- & -- & externally fixed (CEX)\\
Tenure smoothing & 0.000 & -- & -- & -- & externally fixed\\
Child-dependent bequest utility & 0.000 & -- & -- & -- & externally fixed\\
\bottomrule
\end{tabularx}
\end{table}

\end{document}
